\documentclass[11pt]{article}
\usepackage{amsmath, amssymb}
\usepackage{geometry}
\usepackage{array}
\usepackage{booktabs}
\geometry{margin=0.7in}
\setlength{\parindent}{0pt}
\setlength{\parskip}{0.35em}
\renewcommand{\arraystretch}{1.15}

\newcommand{\cheatsheettitle}[2]{%
\begin{center}
{\LARGE #1}\\[0.4em]
{\large #2}
\end{center}
}


\begin{document}
\cheatsheettitle{Algebra I Foundations Cheatsheet}{Module 1}

\section*{Order of Operations}
\[
\text{Parentheses}\rightarrow \text{Exponents}\rightarrow \text{Multiply/Divide}\rightarrow \text{Add/Subtract}
\]
Multiplication and division go left to right. Addition and subtraction go left to right.

\section*{Signed Numbers}
\[
(+)(+)=+,\qquad (-)(-)=+,\qquad (+)(-)=-,\qquad (-)(+)=-
\]
\[
a-b=a+(-b)
\]
Add numbers with the same sign by adding absolute values and keeping the sign. Add numbers with different signs by subtracting absolute values and keeping the sign of the larger absolute value.

\section*{Fractions}
\[
\frac{a}{b}+\frac{c}{d}=\frac{ad+bc}{bd},\qquad
\frac{a}{b}-\frac{c}{d}=\frac{ad-bc}{bd}
\]
\[
\frac{a}{b}\cdot\frac{c}{d}=\frac{ac}{bd},\qquad
\frac{a}{b}\div\frac{c}{d}=\frac{a}{b}\cdot\frac{d}{c}
\]
Restrictions:
\[
b\ne0,\qquad d\ne0
\]

\section*{Evaluating Expressions}
Substitute first, then simplify.
\[
3x^2-2x\quad \text{at }x=-4
\]
means
\[
3(-4)^2-2(-4)
\]
Use parentheses when substituting negative numbers.

\section*{Like Terms}
Like terms have the same variable part and the same exponents.
\[
7x^2-3x^2=4x^2
\]
\[
5a+2b-3a=2a+2b
\]

\section*{Distributive Property}
\[
a(b+c)=ab+ac
\]
\[
-a(b-c)=-ab+ac
\]
Common mistake:
\[
3(x+5)\ne 3x+5
\]

\section*{Basic Exponents and Radicals}
\[
x^1=x,\qquad x^0=1\quad (x\ne0)
\]
\[
x^m\cdot x^n=x^{m+n},\qquad \frac{x^m}{x^n}=x^{m-n}\quad (x\ne0)
\]
\[
(x^m)^n=x^{mn}
\]
\[
\sqrt{x^2}=|x|
\]

\section*{Algebraic Properties}
\[
a+b=b+a,\qquad ab=ba
\]
\[
(a+b)+c=a+(b+c),\qquad (ab)c=a(bc)
\]
\[
a(b+c)=ab+ac
\]

\section*{Quick Checks}
\begin{itemize}
  \item Combine like terms only when the variable parts match exactly.
  \item Keep fraction bars as grouping symbols.
  \item Use parentheses for negative substitutions.
  \item Never divide by zero.
\end{itemize}

\end{document}
